\chapter{量化开发者合规路径：从灰色地带到持牌展业}

本章面向已具备相对稳定盈利策略组合的量化开发者，梳理从「借贷式资管」「代客理财」等非正规安排，过渡到在主要司法辖区合法管理外部资金的典型阶梯。\textbf{本书不构成法律意见}；具体能否展业须由当地持牌律师结合资金性质、募集方式、客户类型与跨境因素认定。

\section{为何「借贷 / 配资 / 代客理财」本身风险高}

中国证监会投资者教育案例明确指出：证券从业人员\textbf{私下代客理财、约定收益分成}，违反《证券法》第143--145条及经纪人相关规则；纠纷中投资者往往只能向个人追责。参见官方案例：\url{http://www.csrc.gov.cn/csrc/c100211/c1452037/content.shtml}。

实务中常与下列风险标签一并讨论（是否构成违法或犯罪须个案认定）：面向\textbf{不特定对象}、公开或变相公开宣传、\textbf{保本保收益}承诺，可能触及非法集资类刑事风险；场外配资长期为监管打击重点——\textbf{结构违法则策略再稳也不安全}。

\section{对外资管的合规阶梯（概念模型）}

\begin{figure}[htbp]
\centering
\begin{tikzpicture}[node distance=1.15cm and 0.9cm]
  \node[compliance=purple] (A) {A. 自有或极窄关系户\\范围与账户出借、税务等仍须个案评估};
  \node[compliance=violet, below=of A] (B) {B. 入职持牌机构\\券商资管、公募、期货资管、银行理财子等};
  \node[compliance=blue!70!cyan, below=of B] (C) {C. 证券投资咨询机构\\证监会许可范围内；咨询 $\neq$ 代客全权委托};
  \node[compliance=teal, below=of C] (D) {D. 私募基金管理人\\中基协登记 + 产品备案（证券类）};
  \node[compliance=green!50!black, below=of D] (E) {E. 投顾输出（如 MOM）\\满足条件后向银行、信托、券商资管提供投顾服务};

  \draw[compliance arrow] (A) -- (B);
  \draw[compliance arrow] (B) -- (C);
  \draw[compliance arrow] (C) -- (D);
  \draw[compliance arrow] (D) -- (E);

  \begin{scope}[on background layer]
    \node[fill=gray!5, rounded corners=14pt, inner sep=14pt, yshift=-2pt,
          fit=(A)(E), draw=gray!40, dashed] {};
  \end{scope}
\end{tikzpicture}
\caption{量化开发者常见合规演进阶梯（示意，非唯一路径）}
\label{fig:quant-ladder}
\end{figure}

对已能稳定盈利的策略组合，实务上常见顺序是：\textbf{B} 验证规模与合规耐受 $\rightarrow$ 有出资人与团队后再上 \textbf{D}；若主要输出「研究/模型」而非直接管账，可重点评估 \textbf{C/E} 与合同边界。

\section{中国大陆：私募登记与系统入口}

根据证监会公开的流程说明，私募基金管理人应向\textbf{基金业协会}办理登记，基金募集完毕后办理备案；系统为 \textbf{AMBERS}：\url{https://ambers.amac.org.cn}。流程索引：\url{http://www.csrc.gov.cn/csrc/c101939/c1045416/content.shtml}。《私募投资基金监督管理暂行办法》第九条明确：协会登记备案\textbf{不构成}对管理人投资能力、持续合规的认可，\textbf{也不保证}基金财产安全。

《私募投资基金登记备案办法》（中基协发〔2023〕5号）及配套指引下，市场常讨论的硬性方向包括（以协会\textbf{最新材料清单}为准）：私募基金管理人实缴货币资本不低于\textbf{1000万元人民币}（或等值可自由兑换货币）等；证券类负责投资的高管常需满足\textbf{最近5年内连续2年以上}担任基金经理或投资决策负责人等，且单只产品管理规模不低于\textbf{2000万元}等证明要求；合规风控负责人须具备合规、风控、法律、会计等相关经验；近年内监管处罚、重大风险机构任职等构成负面情形。

\section{美国（简要锚点）}

管理外部客户资产（含多数对冲基金架构）常需注册为 \textbf{Investment Adviser}，提交 \textbf{Form ADV}，通过 \textbf{IARD}。SEC 与州管辖取决于 AUM 等（\textbf{Rule 203A-1} 等）。入口参考：\url{https://www.sec.gov/divisions/investment/iaregulation/regia.htm}。策略以期货/互换为主时可能叠加 \textbf{CFTC/NFA}。

\section{新加坡（简要锚点）}

从事《证券与期货法》（SFA）项下基金管理，须申请持有 \textbf{CMS} 牌照的 \textbf{LFMC}；仅管理风投基金可申请简化 \textbf{VCFM}。MAS 官方说明含豁免情形、Form 1、eLicensing、审批约6个月、年费 S\$4{,}000 等：\url{https://www.mas.gov.sg/regulation/capital-markets/apply-for-licensing-or-registration-of-capital-market-entities/fund-management-licensing}。

\section{与集团架构、分工条款的衔接}

「母公司—子公司」主要服务治理、合同主体与税务筹划讨论，\textbf{不能替代}管理人登记/产品备案或相应辖区资管牌照。「技术方只写策略、资金方承担一切」在协会材料中，投资或高管岗位仍可能被审核；任职与对外披露须由律师设计。
